\section{Analysis}
\label{sec:analysis}

\paragraph{Task dependence.} Topic labels in AG News hinge on
vocabulary that six layers already capture, and the per-class errors show it.
BERT sends 140 Business stories to Sci/Tech and 90 the other way. DistilBERT
makes 140 and 92. The extra depth finds nothing more in those items. Short SST-2
sentences full of negation and contrast are where the extra layers pay off.

\paragraph{Head versus encoder.} Five heads, from one linear layer
to two hidden layers, fit inside 0.13 points. Fine-tuning shapes the
\texttt{[CLS]} vector until a linear boundary separates the classes, which
leaves extra head layers nothing to do. The linear head
even lifts BERT above its own base run on SST-2 (92.89 against 92.32
accuracy), so the default $\tanh$ pooler adds nothing here.

\paragraph{Latency regimes.} At batch 1 latency stays flat from 64 to 256
tokens (BERT 5.56, 5.61 and 5.11\,ms) because the GPU is waiting on kernel
launches. From batch 8 on, latency tracks the input length. DistilBERT wins
both regimes, with half the kernels to launch in the first and half the
arithmetic in the second. Its memory edge shrinks as activations outgrow the
weights.

\paragraph{Limitations.} All runs use seed 42, so gaps of a few tenths are
ties and we cannot give intervals. Two epochs sit at the low end of what
\citet{devlin2019bert} recommend. Yelp is subsampled for both models alike. Five base runs predate the logging of
validation F1, hence its absence from Table~\ref{tab:results}.
